\section{Reception, Canonizations, and Current Standing}

\subsection{Papal and ecclesial reception}

Reception followed the 1930 judgment in the forms proper to each act,
none of which upgrades the judgment itself: Pius XII's radio messages
and the 1942 and 1952 consecrations (\S6); Paul VI's pilgrimage of
13 May 1967 for the fiftieth anniversary, with the apostolic
exhortation \work{Signum magnum} signed that day; John Paul II's
pilgrimages of 1982, 1991, and 2000; Benedict XVI's of 2010; Francis's
of 2017 and 2023. Benedict XVI's homily at F\'atima on 13 May 2010
gave the reception its lasting formula of continuing relevance: ``We
would be mistaken to think that Fatima's prophetic mission is
complete''---prophetic here in exactly the sense his own 2000
commentary had defined, a summons to conversion, not a further
calendar.\endnote{Paul VI: \work{Signum magnum}, 13 May 1967, Vatican
web text; pilgrimage attested throughout the Sanctuary's records.
Benedict XVI, homily at the Esplanade of the Shrine, 13 May 2010,
English text at
\url{https://www.vatican.va/content/benedict-xvi/en/homilies/2010/documents/hf_ben-xvi_hom_20100513_fatima.html}.
Papal visits establish papal devotion and pastoral government; they
are not additional judgments on the events.}

\subsection{The seers' causes and the centenary canonizations}

Francisco Marto died 4 April 1919, Jacinta 20 February 1920, both of
the pneumonic influenza epidemic. Their joint cause produced decrees
of heroic virtues signed by John Paul II on 13 May 1989---the first
recognition of heroic virtue in non-martyr children, a decision of
general import for the Church's understanding of childhood
sanctity---and he beatified them at F\'atima on 13 May 2000, before
the announcement of the third part's publication. Francis canonized
them at F\'atima on 13 May 2017, during the Mass of the apparitions'
centenary: the first canonized non-martyr children, and the first
saints of the F\'atima event itself. His homily held the canonization
inside the message's own center---``We can take as our examples Saint
Francisco and Saint Jacinta, whom the Virgin Mary introduced into the
immense ocean of God's light and taught to adore him''---and quoted
the Third Memoir's account of Jacinta's vision of a suffering
pope.\endnote{Deaths and 1989 decrees: Sanctuary seer chronology.
Beatification: John Paul II, homily at F\'atima, 13 May 2000, Vatican
web text. Canonization: Francis, homily at the Holy Mass with the
rite of canonization of Blesseds Francisco Marto and Jacinta Marto,
Shrine of Our Lady of F\'atima, 13 May 2017, English text at
\url{https://www.vatican.va/content/francesco/en/homilies/2017/documents/papa-francesco_20170513_omelia-pellegrinaggio-fatima.html},
quoting \work{Memoirs} III, 6. The Latin formula of canonization
pronounced in the rite is in the celebration booklet, not
re-verified here. A canonization certifies the sanctity of the
persons; it does not adjudicate the apparitions' every detail.}
Sister Lucia died in the Carmel of Coimbra on 13 February 2005;
Benedict XVI dispensed from the customary five-year wait so that her
cause could open in Coimbra in 2008, and Pope
Francis approved the decree of her heroic virtues on 22 June 2023,
giving her the title Venerable; her beatification awaits an approved
miracle.\endnote{Dicastery for the Causes of Saints, decree on the
heroic virtues of the Servant of God Maria Lucia of Jesus and of the
Immaculate Heart, 22 June 2023 (Portuguese text in the Sanctuary's
documentation collection; References); waiver and diocesan phase:
Sanctuary seer chronology.}

\subsection{The 2016 communiqu\'e on the ``incomplete publication'' claims}

In 2016 several outlets circulated statements attributed to the
German professor Ingo Dollinger, according to which Cardinal
Ratzinger had privately confided, after June 2000, that the
publication of the third secret was incomplete. The Holy See Press
Office answered on 21 May 2016 with a communiqu\'e whose operative
sentences deserve exact quotation: ``Pope emeritus Benedict XVI
declares `never to have spoken with Professor Dollinger about
Fatima', clearly affirming that the remarks attributed to Professor
Dollinger on the matter `are pure inventions, absolutely untrue', and
he confirms decisively that `the publication of the Third Secret of
Fatima is complete'.''\endnote{Holy See Press Office,
``Communiqu\'e on articles regarding the Third Secret of Fatima,''
21 May 2016, English bulletin text at
\url{https://press.vatican.va/content/salastampa/en/bollettino/pubblico/2016/05/21/160521b.html}.
The claims answered are identified in the communiqu\'e's own first
paragraph; they are recited here only as the object of the Holy
See's response.} With that declaration, every layer of the textual
question---manuscript, facsimile, translation, completeness---stands
publicly closed by the persons with direct knowledge.

\subsection{Standing under the 2024 norms}

The Dicastery for the Doctrine of the Faith's \work{Norms for
Proceeding in the Discernment of Alleged Supernatural Phenomena}
(approved 4 May 2024; effective 19 May 2024, Pentecost) replaced the
1978 norms and restructured conclusions into six determinations:
\term{Nihil obstat}; \term{Prae oculis habeatur}; \term{Curatur};
\term{Sub mandato}; \term{Prohibetur et obstruatur}; and
\term{Declaratio de non supernaturalitate}. As a rule, ``neither the
Diocesan Bishop, nor the Episcopal Conferences, nor the Dicastery
will declare that these phenomena are of supernatural origin, even if
a \term{Nihil obstat} is granted,'' though ``the Holy Father can
authorize a special procedure in this regard.''\endnote{DDF,
\work{Norms}, 17 May 2024, English text at
\url{https://www.vatican.va/roman_curia/congregations/cfaith/documents/rc_ddf_doc_20240517_norme-fenomeni-soprannaturali_en.html}:
conclusions at nn.~17--22, the rule at n.~23; approval, effective
date, and replacement of the 25 February 1978 norms in the final
formula and art.~27.} Three consequences for F\'atima, stated
carefully. First, the norms govern \emph{proceedings}; they do not
retroactively reclassify legacy decisions, and the 1930 declaration
remains what it always was---a positive diocesan judgment in its own
historical formula (``worthy of belief''), which this study does not
relabel as a 2024-vocabulary \term{Nihil obstat}. Second, the norms'
own presentation cites F\'atima as the standing example of what
approval means, quoting the 2000 explanation that ``ecclesiastical
approval of a private revelation highlights that the message contains
nothing contrary to faith or morals'' (26 June 2000). Third, no act
of the Dicastery under the new norms concerns F\'atima: the DDF's
public documents index contains, for F\'atima, only the 2000
\work{Message of Fatima} dossier (checked 2026-07-25), and no later
Holy See intervention superseding the 1930 judgment or the 2000
interpretation was located in the sources searched. Silence is
reported as a bounded negative search, not as a guarantee.\endnote{%
DDF documents index at \url{https://www.doctrinafidei.va/en/documenti.html},
listing ``Documents Regarding The Message of Fatima (26 June 2000)''
as the Dicastery's F\'atima entry; search bounds in Appendix~A.}

\begin{statusdossier}{Controlling-status dossier --- F\'atima, as of
2026-07-25}
\dosfield{Event identity}{Six reported apparitions of Our Lady to
Lucia dos Santos, Francisco Marto, and Jacinta Marto at the Cova da
Iria, parish of F\'atima, 13 May--13 October 1917 (August encounter
at Valinhos, 19 August), with the reported solar phenomenon of
13 October 1917}
\dosfield{Jurisdiction}{Diocese of Leiria (now Leiria-F\'atima),
Latin Church, Portugal}
\dosfield{Competent act}{Carta Pastoral of Bishop Jos\'e Alves
Correia da Silva, 13 October 1930, Portuguese; witness:
\work{Documenta\c{c}\~ao Cr\'itica de F\'atima: Sele\c{c}\~ao}
(2013), Doc.~133}
\dosfield{Procedural regime}{Pre-1978 diocesan-ordinary discernment;
process opened 3 May 1922; commission report 13 April 1930}
\dosfield{Operative formula}{``declarar como dignas de cr\'edito as
vis\~oes das crian\c{c}as\ldots\ nos dias 13 de maio a outubro de
1917'' and ``permitir oficialmente o culto de Nossa Senhora de
F\'atima''}
\dosfield{Exact object}{The children's 1917 Cova da Iria visions and
the public cult; not the 1916 angel cycle, the 1915 reports, the 1921
experience, Pontevedra 1925--1926, Tuy 1929, the memoirs, or the
secret's texts}
\dosfield{Later acts}{Custody and publication of the third part (CDF,
26 June 2000, with theological commentary); papal consecrations 1942,
1952, 1984, 2022; beatification 2000 and canonization 2017 of
Francisco and Jacinta; Lucia declared Venerable 2023; 2016
communiqu\'e affirming complete publication}
\dosfield{Current status}{Positive legacy diocesan judgment in force;
2000 CDF interpretation controlling for the secret; no DDF act under
the 2024 norms; no located superseding or restricting act}
\dosfield{What it does not establish}{Obligation of faith; verbal
inerrancy of reported words (see the October war statement); a
defined mechanism for the solar phenomenon; authentication of every
later text, promise, or private interpretation}
\end{statusdossier}
